\section{Zuverlässigkeit}

Zuverlässigkeit bezeichnet die Wahrscheinlichkeit, mit der ein System über einen bestimmten Zeitraum hinweg erreichbar ist.
Dabei besteht Zuverlässigkeit aber nicht nur aus der Verfügbarkeit eines Systems, sondern auch aus dessen Fehlertoleranz, Stabilität und Wiederherstellbarkeit.

Die Fehlertoleranz beschreibt dabei nicht, dass keine Fehler passieren, sondern wie ein System mit Fehlern umgeht. Ein zuverlässiges System soll trotz Fehler weiterlaufen, ohne dass es zu Ausfällen kommt oder sich Fehler anderweitig bemerkbar machen.
Daher ist es bei der Zuverlässigkeit nicht am wichtigsten, dass ein System keine Fehler produziert, sondern dass es kontrolliert mit Fehlern umgehen kann.

Wenn es zu Ausfällen im System kommt – durch Fehler oder auch Updates –, muss das System automatisch darauf reagieren können, damit die Downtime möglichst gering ist.
Daher muss das System nach einem Ausfall wieder in seinen ursprünglichen, funktionierenden Zustand versetzt werden.

Zur Fehlervermeidung ist es wichtig, dass ein System stabil läuft, damit es funktionsfähig bleibt. 
Dazu gehört, dass ein System mit Last umgehen kann und sich entsprechend der Last hochskalieren kann, um eine Überlastung des Systems zu vermeiden.

Wenn ein System eine schlechte Performance hat, kann die Zuverlässigkeit geringer sein. Dies hängt damit zusammen, dass die Stabilität des Systems gefährdet wird, da mehr Last entsteht.
